\documentclass[11pt]{article}
\usepackage[T1]{fontenc}
\usepackage[utf8]{inputenc}
\usepackage[margin=1in]{geometry}
\usepackage{amsmath,amssymb,amsthm,booktabs}
\usepackage[hidelinks]{hyperref}
\hypersetup{
 pdftitle={An even-strand Markov calculus for classical and virtual links},
 pdfauthor={Alec Kriebel}
}
\newtheorem{theorem}{Theorem}
\newtheorem{corollary}[theorem]{Corollary}
\newcommand{\sig}{\sigma}
\newcommand{\arr}{\longleftrightarrow}
\newcommand{\WB}[2]{W_{#1}(#2)}
\title{An even-strand Markov calculus for classical and virtual links}
\author{Alec Kriebel\\\small ORCID: \href{https://orcid.org/0009-0001-9320-500X}{0009-0001-9320-500X}}
\date{}
\begin{document}
\maketitle
\begin{abstract}
We give an explicit affirmative answer to the ordinary-closure interpretation of
Problem~42 in the 2014 survey of Fenn, Ilyutko, Kauffman and Manturov.
Four reversible algebraic schemes classify oriented classical link closures
using only even-strand braid states. Four additional exchange schemes give the
corresponding statement for oriented virtual links. The proof is an elementary
conversion of the established unrestricted Markov theorems: positively stabilize
each odd-strand vertex and lift every generating edge. Arbitrary finite word
blocks are allowed; minimality and uniform geometric locality are not asserted.
No historical-priority or present-openness claim is made.
\end{abstract}

\section{Scope and conventions}
Problem~42 of~\cite[p.~37]{FIKM}, also on p.~34 of its author preprint,
asks for a Markov formulation involving only even-strand braids.
We retain \emph{ordinary oriented, unframed closure} throughout. We supply both
classical and virtual versions, since the item's braid terminology is unqualified.
Plat closure, framing and transverse isotopy are different classifications.
An even-strand plat theorem would not answer the ordinary-closure statement:
the identity two-strand braid has a two-component ordinary closure and a
one-component plat closure.

A state is a tagged word $(N,w)$, where $N\ge2$ is even and all generator
indices lie in $\{1,\ldots,N-1\}$. Tags cannot be suppressed: empty words on
two and four strands close to different unlinks. Words multiply by concatenation;
$w^{-1}$ reverses the letters and inverts each one. Classical words use
$\sig_i^{\pm1}$; virtual words additionally use $v_i=v_i^{-1}$.
Write $\WB{N}{k}$ for words whose indices are at most $k$, always regarded as
words on the stated even $N$ strands. The empty word is allowed, also for $k=0$.
These are syntactic support conditions, not tests of membership up to an
unspecified braid equivalence. Let $s(w)$ increase every index by one, retaining
the types and signs of letters. Put
\[
 G_i^{\mathrm{cl}}=\{\sig_i,\sig_i^{-1}\},\qquad
 G_i^{\mathrm{vir}}=\{\sig_i,\sig_i^{-1},v_i\}.
\]
Use the appropriate $G_i$ below.

At every fixed even level we allow inverse cancellation and the Artin relations
\[
 \sig_i\sig_j=\sig_j\sig_i\ (|i-j|>1),\qquad
 \sig_i\sig_{i+1}\sig_i=\sig_{i+1}\sig_i\sig_{i+1}.
\]
For virtual words we additionally allow exactly
\begin{align*}
 v_i^2&=1,&v_i v_j&=v_j v_i\quad(|i-j|>1),&
 v_i v_{i+1}v_i&=v_{i+1}v_i v_{i+1},\\
 \sig_i v_j&=v_j\sig_i\quad(|i-j|>1),&&
 \sig_i v_{i+1}v_i&=v_{i+1}v_i\sig_{i+1}.
\end{align*}
All these relations may be used in arbitrary word contexts, in either direction.
No welded forbidden relation is imposed.

\section{The even-endpoint schemes}
All following arrows are reversible and have only even-strand endpoints.
The four common schemes are
\begin{align}
 (N,b)&\arr(N,aba^{-1}), && a,b\in\WB{N}{N-1}, \tag{C}\label{C}\\
 (N,b\sig_{N-1})&\arr(N,aba^{-1}\sig_{N-1}),
 &&a,b\in\WB{N}{N-2}, \tag{BC}\label{BC}\\
 (N,b\sig_{N-1})&\arr(N,bg),
 &&b\in\WB{N}{N-2},\ g\in G_{N-1}, \tag{T}\label{T}\\
 (N,b)&\arr(N+2,bg\sig_{N+1}),
 &&b\in\WB{N}{N-1},\ g\in G_N. \tag{D}\label{D}
\end{align}
In (D), the old word is placed on the first $N$ strands of the new $N+2$-strand
state. The identity choice in (T) may be omitted. These four schemes,
with classical letters only, form the classical calculus.

The virtual calculus includes four further schemes. For
$a,b\in\WB{N}{N-2}$, allow
\begin{align}
 (N,a\sig_{N-1}^{-1}b\sig_{N-1})
 &\arr(N,av_{N-1}bv_{N-1}), \tag{R}\label{R}\\
 (N,s(a)\sig_1^{-1}s(b)\sig_1)
 &\arr(N,s(a)v_1s(b)v_1). \tag{L}\label{L}
\end{align}
For even $N\ge4$ and $a,b\in\WB{N}{N-3}$, allow
\begin{align}
 (N,a\sig_{N-2}^{-1}b\sig_{N-2}\sig_{N-1})
 &\arr(N,av_{N-2}bv_{N-2}\sig_{N-1}), \tag{BR}\label{BR}\\
 (N,s(a)\sig_1^{-1}s(b)\sig_1\sig_{N-1})
 &\arr(N,s(a)v_1s(b)v_1\sig_{N-1}). \tag{BL}\label{BL}
\end{align}
In (BL), the shifted blocks have indices at most $N-2$; the separate retained
terminal crossing has index $N-1$ and is not shifted.
There are four classical or eight virtual \emph{scheme families}, apart from the
defining relations. Blocks may have arbitrary finite length. Neither a minimal
list nor a bound on the number of strands in a geometric support disk is claimed.

\begin{theorem}\label{main}
Every nonempty oriented classical link is the ordinary closure of an even-strand
classical braid. Two such tagged braids have isotopic oriented closures if and
only if they are connected by defining relations at even levels and the schemes
(C), (BC), (T), (D), with classical letters only.

Every nonempty oriented virtual link has an even-strand virtual braid
representative. Two such tagged virtual braids have equivalent oriented virtual
closures if and only if they are connected by defining virtual braid relations
at even levels and (C), (BC), (T), (D), (R), (L), (BR), (BL).
\end{theorem}

\section{Proof by positive padding}
The classical inputs are Alexander's representation theorem and ordinary
Markov equivalence by conjugation and signed right stabilization; see
\cite[Theorems~2.1--2.2]{GKS}. The virtual inputs are Kamada's
Proposition~3.1 and Theorem~3.2 in~\cite{Kamada}. In addition to relations and
conjugation, they permit right stabilization by $\sig_m^{\pm1}$ or $v_m$,
and both exchanges at total count $m$:
\begin{equation}\label{unrestricted-exchange}
 a\sig_{m-1}^{-1}b\sig_{m-1}\arr av_{m-1}bv_{m-1},\qquad
 s(a)\sig_1^{-1}s(b)\sig_1\arr s(a)v_1s(b)v_1,
\end{equation}
where the unshifted blocks use indices at most $m-2$.
The signs and shift are also stated in~\cite[\S5, p.~30]{KL}.
These unrestricted theorems are imported, not reproved here.

\begin{proof}[Proof of Theorem~\ref{main}]
First verify soundness without using completeness. (C) is conjugation.
For (BC), the two prefixes on $N-1$ strands are conjugate, and both endpoints
are their positive right stabilizations. (T) compares two permitted right
stabilizations of the same prefix on $N-1$ strands. (D) is a permitted
stabilization followed by a positive stabilization. (R) and (L) are
\eqref{unrestricted-exchange}; in (BR) and (BL), apply that exchange on
$N-1$ strands and then positively stabilize both sides. The support bounds make
each auxiliary operation legitimate. Thus every listed replacement preserves
ordinary oriented closure. Odd-strand objects used to verify this fact are
auxiliary proof objects, not allowed states of the displayed calculus.

For completeness, define, on unrestricted tagged words with $m\ge1$,
\begin{equation}\label{padding}
 P(m,w)=
 \begin{cases}
 (m,w),&m\text{ even},\\
 (m+1,w\sig_m),&m\text{ odd}.
 \end{cases}
\end{equation}
Its closure agrees with the original closure by positive stabilization, and it
fixes every even tagged word exactly. Given an unrestricted Markov chain between
even endpoints, apply $P$ to every vertex. We check its edges exhaustively.

A defining relation at even count remains unchanged. At odd count the same
relation is applied inside its entire old word context, with the positive tail
retained. Right inclusion respects every defining relation; hence no group-word
ambiguity is hidden by padding.

Conjugation $b\arr aba^{-1}$ at even count $m$ becomes (C).
At odd $m$ it becomes (BC) at $N=m+1$, since both old blocks have indices
at most $m-1=N-2$.
For stabilization $(m,b)\arr(m+1,bg)$, $g\in G_m$, the two images are
\[
 \begin{cases}
 (m,b)\arr(m+2,bg\sig_{m+1}),&m\text{ even}:\ (D),\\
 (m+1,b\sig_m)\arr(m+1,bg),&m\text{ odd}:\ (T).
 \end{cases}
\]
This includes both classical signs and virtual stabilization. Reversed arrows
also cover every destabilization.

For a right exchange at even total count $m$, the images give (R) at $N=m$.
At odd total count $m$, its exchange index is $m-1=N-2$, its unshifted blocks
have indices at most $m-2=N-3$, and padding appends $\sig_m=\sig_{N-1}$:
this is (BR). A left exchange similarly gives (L) at even $m$ and (BL) at odd
$m$. In the latter, only the old blocks are shifted, so their indices are at
most $N-2$ and the padding crossing remains separate. The first possible odd
exchange count is $m=3$, explaining $N\ge4$ for (BR) and (BL).

This covers every unrestricted edge. Delete identical successive images if
desired. The resulting chain has only even states and the original endpoints.
The classical chain needs no virtual letters or exchanges. Soundness proves the
converse, and the corresponding Alexander theorem followed by $P$ proves the
representation statements.
\end{proof}

\begin{corollary}
An unrestricted certificate with maximum strand count $M$ converts to an even
certificate with maximum strand count $2\lceil M/2\rceil$.
\end{corollary}
\begin{proof}
Each vertex count $m$ becomes $m+(m\bmod2)$, and every lifted edge above has
exactly its displayed endpoints. This function is nondecreasing in $m$.
\end{proof}
This is a bound for a \emph{supplied} certificate, not an algorithm or complexity
bound for finding one. At one strand the empty word pads to $(2,\sig_1)$,
the unknot. At $N=2$, (BC)'s blocks are empty and the virtual exchanges are
identities. If the empty link is desired, adjoin an isolated zero-strand empty
state; no positive-strand ordinary closure is empty.

\section{Double moves and checkable support boundaries}
Fiedler~\cite{Fiedler} reports that conjugation and double Markov moves alone
do not replace ordinary Markov equivalence. Our calculus includes the additional
(BC) and (T) schemes. Indeed, (D) followed by (T) at $N+2$ gives every classical
double stabilization
\[
 (N,b)\arr (N+2,b\sig_N^\varepsilon\sig_{N+1}^\delta),
 \qquad \varepsilon,\delta\in\{+1,-1\},
\]
because the prefix $b\sig_N^\varepsilon$ has indices at most $N$.
This does not show that (BC) or (T) is generated by conjugation and double moves.
Our comparison uses Fiedler's publisher abstract and identified publisher
first-page preview; we do not independently certify his full counterexample
proof. The elementary padding derivation establishes the stated formulation
without a historical-priority claim or an assertion that the 2014 question
remains open today.

An earlier related announcement is Nencka's two-page contribution
\cite{Nencka}. After introducing a new representation by interwoved strings,
it states a ``Markov generalized equivalence'' theorem using conjugacy at
$2^n$ strands and signed stabilization tails from $2^n$ to $2^{n+1}$.
For $n\ge1$ these are even counts; each such tail is a composite of double
stabilizations by grouping its $2^n$ crossings into successive pairs.
Thus a proved ordinary-closure interpretation of that announcement would be
materially relevant to the classical question. The contribution does not
define its generalized equivalence or supply a proof sufficient to verify
that interpretation, and its convention for possible one-strand participation
is not specified. We neither certify an earlier ordinary-closure resolution
nor refute its announced result. The current proof uses the explicitly stated
ordinary-closure inputs above and includes (BC) and (T); the announcement
supplies no virtual exchange calculus. This distinction does not establish
historical priority for our formulation.
We have identified fuller related texts~\cite{Nencka98,Nencka99}, whose bodies
have not yet been accessed; the bearing of those texts on exact
ordinary-closure priority remains unresolved.

The block restrictions have mathematical content. If (T) at $N=2$ were allowed
with $b=\sig_1^2$, it would relate $\sig_1^3$ to $\sig_1$ after cancellation.
Their closures both have one component, but their Fox $3$-coloring counts are
$9$ and $3$. With the positive generator drawn with its left strand over the right,
the color map is $R(x,y)=(2x-y,x)$ over $\mathbb F_3$; $R^3=1$ and $R$ fixes precisely $x=y$.
The invertible coloring rule gives bijections under the Reidemeister moves,
so these counts rule out that enlarged scheme. The true (T) requires an empty
prefix here.

There is a virtual obstruction missed by component counts as well. Enlarging
(R) at $N=2$ to allow $a=b=\sig_1$ would relate $\sig_1^2$ to
$\sig_1v_1\sig_1v_1$. Their signed ordered intercomponent crossing matrices are
respectively $\left(\begin{smallmatrix}0&1\\1&0\end{smallmatrix}\right)$ and
$\left(\begin{smallmatrix}0&2\\0&0\end{smallmatrix}\right)$, distinct even after
simultaneous component relabeling. These matrices are invariant under oriented
virtual Reidemeister moves: RI adds only self-crossings, classical RII cancels
signs, RIII preserves each ordered component pair, and virtual detours introduce
no classical crossings.
The genuine (R) excludes these blocks.

The accompanying standard-library Python diagnostics check explicit lifted
endpoints, every relation family in bounded word contexts, support and tags,
and exact countercontrols such as these. Finite diagnostics supplement the
universal proof. Literal nonempty left-exchange endpoints transcribed from the
primary formulas check the index shift and separate padding tail without using
the shared shift helper to construct expected endpoints. Zero-shift
countercontrols are separated by ordered crossing matrices despite equal
component counts. These diagnostics are not a link-equivalence oracle, a braid word-problem
solver or formal verification of the imported theorems.

\paragraph{AI and review disclosure.}
AI tools were used extensively in constructing and drafting this formulation,
and in adversarial review and computational verification. This note is
unrefereed; no human peer review or formal proof certification is claimed.

\begingroup
\small
\begin{thebibliography}{9}
\bibitem{FIKM} R.~Fenn, D.~P.~Ilyutko, L.~H.~Kauffman and V.~O.~Manturov,
\emph{Unsolved problems in virtual knot theory and combinatorial knot theory},
Banach Center Publications \textbf{103} (2014), 9--61.
\href{https://doi.org/10.4064/bc103-0-1}{doi:10.4064/bc103-0-1};
\href{https://arxiv.org/abs/1409.2823v1}{arXiv:1409.2823v1}.
\bibitem{GKS} E.~Gorsky, O.~Kivinen and J.~Simental,
\emph{Algebra and geometry of link homology: Lecture notes from the IHES 2021 Summer School},
Bull. London Math. Soc. \textbf{55} (2023), no.~2, 537--591.
\href{https://doi.org/10.1112/blms.12761}{doi:10.1112/blms.12761};
\href{https://arxiv.org/abs/2108.10356v1}{arXiv:2108.10356v1}.
\bibitem{Kamada} S.~Kamada,
\emph{Braid presentation of virtual knots and welded knots},
\href{https://arxiv.org/abs/math/0008092v1}{arXiv:math/0008092v1} (2000),
\S\S2--3, Proposition~3.1 and Theorem~3.2.
\bibitem{KL} L.~H.~Kauffman and S.~Lambropoulou,
\emph{Virtual braids and the $L$-move},
J. Knot Theory Ramifications \textbf{15} (2006), 773--811;
\href{https://arxiv.org/abs/math/0507035v3}{arXiv:math/0507035v3}, \S5, p.~30.
\bibitem{Nencka} H.~Nencka,
\emph{Generalization of the Markov theorem and Cantorian-like braid groups},
in \emph{Topology and Applications: International Topological Conference
dedicated to P.~S.~Alexandroff's 100th Birthday, Moscow, May 27--31, 1996},
PHASIS, Moscow (1996), 147--148.
\href{https://ru.djvu.online/file/OzKCMNnU4ojPn}{ISBN 5-7036-0017-0; scanned contribution}.
\bibitem{Nencka98} H.~Nencka, \emph{Cantorian braid groups},
Methods Funct. Anal. Topology \textbf{4} (1998), no.~2, 66--75.
\href{https://mfat.imath.kiev.ua/article/?id=71}{Publisher record; MR1770817}.
\bibitem{Nencka99} H.~Nencka, \emph{On some extensions of Artin's braid relations},
Contemporary Mathematics \textbf{233} (1999), 221--233.
\href{https://doi.org/10.1090/conm/233/03432}{doi:10.1090/conm/233/03432}.
\bibitem{Fiedler} T.~Fiedler,
\emph{Markov moves cannot be replaced by double Markov moves},
J. Knot Theory Ramifications \textbf{12} (2003), no.~4, 575--577.
\href{https://doi.org/10.1142/S0218216503002561}{doi:10.1142/S0218216503002561}.
\end{thebibliography}
\endgroup
\end{document}
